% Living Showcase: optionale Code-Module (Palette Slot 7).
\StartChapter[ch:code]{Code und Kommentare}

\SubsectionBar[sec:code-listing]{CodeExpert-Listing}

\begin{codebox}[Was der Listing-Stil einfärbt\ZSFTitleTag{codebox}]
\begin{lstlisting}[style=CodeExpert]
# Kommentar: gruen gesetzt
def hervorhebung(text="Zeichenkette", zahl=42):
    if zahl > 0:
        return text
    return None
\end{lstlisting}
\end{codebox}

\SubsectionBar[sec:code-split]{Teilbare Code-Gruppe}

\begin{codebox}[Teil 1 trägt den Titel\ZSFTitleTag{part=first}][part=first]
\begin{lstlisting}[style=CodeExpert]
# Rahmen oben, unten offen
werte = [1, 2, 3]
\end{lstlisting}
\end{codebox}
\begin{codebox}[][part=mid]
\begin{lstlisting}[style=CodeExpert]
# part=mid: beide Kanten offen
summe = sum(werte)
\end{lstlisting}
\end{codebox}
\begin{codebox}[][part=last]
\begin{lstlisting}[style=CodeExpert]
# part=last: Rahmen unten
print(summe)
\end{lstlisting}
\end{codebox}

\SubsectionBar[sec:code-comments]{Smarte Zeilenkommentare}

\begin{statementbox}[Kommentar rechts oder darüber\ZSFTitleTag{CodeLine}]
  \ttfamily
  \CodeLine{kurz = 1}[passt rechts daneben]\\
  \CodeLine{sehr_langer_variablenname = 1}[zu breit, deshalb darüber]\\
  \CodeLine{limit = 10}\InlineComment{InlineComment: erzwungen rechts}\\
  \OverlineComment{OverlineComment: erzwungen darüber}
  \CodeLine{fertig = True}
\end{statementbox}

\begin{codebox}
\begin{lstlisting}[style=CodeExpert]
# codebox ohne Titelargument: keine Kopfzeile
titellos = True
\end{lstlisting}
\end{codebox}

\begin{factlist}
  \ZSFFact{\ZSFkeyword{CodeLine}}{Entscheidet nach Breite, wo der Kommentar steht.}
  \ZSFFact{SetCodeCommentThreshold}{Verschiebt die Schwelle für den Umbruch.}
  \ZSFFact{ResetCodeCommentThreshold}{Stellt den Standard wieder her.}
\end{factlist}
